% Zuordnungen – bank/zuordnungen/ (zusammenbau v0.4, Heft)
\einheitenkopf[t6]{Zuordnungen}
\setcounter{aufgabe}{46}

% Anlauf (ohne Prüfkennung)
\begin{aufgabe}{Darstellen – Anlauf}
\begin{teile}
\teil Ein Radfahrer fährt gleichmäßig: nach 1 Stunde 15 km, nach 2 Stunden 30 km, nach 3 Stunden 45 km. Trage die Wertepaare in die Tabelle ein.

\wertetabelleleer{Zeit in h}{Weg in km}{3}
\teil Trage die Punkte aus der Tabelle ins Koordinatensystem ein.

\wertetabelle[2,4,6,8]{Zeit in h}{Weg in km}{1,2,3,4} \begin{ksys}[xmin=0,xmax=5,ymin=0,ymax=9,xlabel=Zeit in h,ylabel=Weg in km]\end{ksys}
\teil Mia hat die Tabelle ins Koordinatensystem übertragen (Bild). Finde den Fehler.

\wertetabelle[3,6]{Zeit in h}{Weg in km}{1,2} \begin{ksys}[xmin=0,xmax=7,ymin=0,ymax=7,xlabel=Zeit in h,ylabel=Weg in km]\punkt{3}{1}{A}\punkt{6}{2}{B}\end{ksys}
\end{teile}
\end{aufgabe}

\begin{aufgabe}{Darstellen – Prüfungsaufgaben}
\begin{teile}
\teil Ein Sportverein zählt die Abonnenten seines Kanals. Trage die Werte aus der Tabelle ins Koordinatensystem ein. Die senkrechte Achse ist in Schritten von 50 beschriftet, ein Kästchen steht für 10. \hfill (P10 2025 OS)

\wertetabelle[70,95,135,185,260]{Zeit in Monaten}{Abonnenten}{0,1,2,3,4} \begin{ksys}[xmin=0,xmax=5,ymin=0,ymax=300,ystep=50,xlabel=Zeit in Monaten,ylabel=Abonnenten]\end{ksys}
\teil Auf einem Teich werden die Seerosenblätter gezählt. Trage die Werte aus der Tabelle ins Koordinatensystem ein. Die senkrechte Achse ist in Schritten von 50 beschriftet, ein Kästchen steht für 10. \hfill (P10 2025 OS)

\wertetabelle[40,60,85,125,180]{Zeit in Wochen}{Blätter}{0,1,2,3,4} \begin{ksys}[xmin=0,xmax=5,ymin=0,ymax=200,ystep=50,xlabel=Zeit in Wochen,ylabel=Blätter]\end{ksys}
\end{teile}
\end{aufgabe}

% Anlauf (ohne Prüfkennung)
\begin{aufgabe}{Achseneinteilung – Anlauf}
\begin{teile}
\teil Auf einer Achse stehen die Zahlen 0, 20 und 40, je 4 Kästchen auseinander. Wie viel ist ein Kästchen wert? 1 Kästchen = \leerfeld
\teil Die Werte gehen bis 60 km, die Achse hat 15 Kästchen. Welche Einteilung passt? \\ \kreuz{1 Kästchen = 1 km} \\ \kreuz{1 Kästchen = 5 km}
\teil Wähle für die waagerechte Achse eine Einteilung und zeichne die beschriftete Achse.

\wertetabelle[2,4,6,8,10,12]{Anzahl Hefte}{Preis in €}{1,2,3,4,5,6}
\end{teile}
\end{aufgabe}

\begin{aufgabe}{Achseneinteilung – Prüfungsaufgaben}
\begin{teile}
\teil Ein Regenfass enthält anfangs 90 l Wasser und läuft gleichmäßig aus. Ergänze die Tabelle. Zeichne dann auf Karopapier ein Koordinatensystem mit 20 Kästchen nach rechts und 20 Kästchen nach oben: Wähle eine Einteilung, beschrifte beide Achsen und trage die Punkte ein. \hfill (P10 2021 OS)

\wertetabelle[,80,70,60,50,40,]{Zeit in min}{Inhalt in l}{0,5,10,15,20,25,30}
\teil Ein Akku wird gleichmäßig geladen. Ergänze die Tabelle. Zeichne dann auf Karopapier ein Koordinatensystem mit 26 Kästchen nach rechts und 20 Kästchen nach oben: Wähle eine Einteilung, beschrifte beide Achsen und trage die Punkte ein. \hfill (P10 2021 OS)

\wertetabelle[,25,35,45,55,65,]{Zeit in min}{Ladung in \%}{0,10,20,30,40,50,60}
\end{teile}
\end{aufgabe}

% Anlauf (ohne Prüfkennung)
\begin{aufgabe}{Antiproportional – Anlauf}
\begin{teile}
\teil Je mehr Maler an einer Wand arbeiten, desto … Zeit brauchen sie. \\ \kreuz{mehr} \\ \kreuz{weniger}
\teil 2 Bagger brauchen für eine Baugrube 10 Tage. Wie lange brauchen 4 Bagger? \leerfeld[Tage]
\teil 2 Traktoren pflügen ein Feld in 9 Stunden. Wie lange brauchen 6 Traktoren? \leerfeld[h]
\end{teile}
\end{aufgabe}

\begin{aufgabe}{Antiproportional – Prüfungsaufgaben}
\begin{teile}
\teil Ein Futtervorrat reicht für 6 Ziegen 10 Tage. Wie viele Tage reicht er für 5 Ziegen? \hfill (P10 2014 GYM) \\ \leerfeld[Tage]
\teil Ein Wasservorrat reicht für 3 Wanderer 10 Tage. Wie viele Tage reicht er für 5 Wanderer? \hfill (P10 2014 GYM) \\ \leerfeld[Tage]
\end{teile}
\end{aufgabe}

% Anlauf (ohne Prüfkennung)
\begin{aufgabe}{Rate – Anlauf}
\begin{teile}
\teil Auf dem Tankbeleg steht: 40 l, 1,79 € je Liter, gesamt 71,60 €. Welche Angabe ist die Rate?
\teil Äpfel kosten 2,50 € je kg. Was kosten 4 kg? \leerfeld[€]
\teil Kartoffeln kosten 1,50 € je kg. Wie viele Kilogramm bekommt man für 6 €? \leerfeld[kg]
\end{teile}
\end{aufgabe}

\begin{aufgabe}{Rate – Prüfungsaufgaben}
\begin{teile}
\teil Ein Lieferwagen verbraucht 8 l Diesel auf 100 km und fährt im Jahr 30\,000 km. Ein Liter Diesel kostet 162,5 ct. Berechne die Dieselkosten für ein Jahr in Euro. \hfill (P10 2014 OS) \\ \leerfeld[€]
\teil Ein Traktor verbraucht 15 l Diesel je Arbeitsstunde und läuft im Jahr 400 Stunden. Ein Liter Diesel kostet 148 ct. Berechne die Dieselkosten für ein Jahr in Euro. \hfill (P10 2014 OS) \\ \leerfeld[€]
\teil Ein Taxi fährt im Jahr 50\,000 km und verbraucht 7 l Diesel auf 100 km. Im letzten Jahr kostete ein Liter 150 ct, in diesem Jahr 158 ct. Der Fahrer sagt: „Bei gleicher Fahrleistung zahle ich 280 € mehr.“ Weise nach, dass das stimmt. \hfill (P10 2014 OS)
\teil Eine Schule verbraucht an jedem ihrer 200 Heiztage 90 l Heizöl. Der Preis stieg von 95 ct auf 1,05 € je Liter. Der Hausmeister sagt: „Wir zahlen jetzt 1\,800 € mehr.“ Weise nach, dass das stimmt. \hfill (P10 2014 OS)
\teil Eine Regentonne ersetzt im Garten 25 l Leitungswasser am Tag. 1 l Leitungswasser kostet 0,4 ct. Wie viel Geld spart man am Tag? Gib das Ergebnis in Euro an. \hfill (P10 2024 OS) \\ \leerfeld[€]
\teil Ein Sparduschkopf spart 30 l Warmwasser am Tag. 1 l Warmwasser kostet 1,2 ct. Wie viel Geld spart man am Tag? Gib das Ergebnis in Euro an. \hfill (P10 2024 OS) \\ \leerfeld[€]
\steil Ein Aufzug in einem Aussichtsturm ist 10 s nach dem Start in 150 m Höhe und nach weiteren 40 s in 310 m Höhe. Er fährt gleichmäßig. Berechne seine Geschwindigkeit in km/h. \hfill (P10 2015 OS) \\ \leerfeld km/h
\steil Ein Heißluftballon sinkt gleichmäßig. 30 s nach Beginn des Sinkflugs ist er in 800 m Höhe, nach weiteren 200 s in 400 m Höhe. Berechne seine Sinkgeschwindigkeit in km/h. \hfill (P10 2015 OS) \\ \leerfeld km/h
\end{teile}
\end{aufgabe}

\begin{aufgabe}{Rate – Prüfungsaufgaben – weiter}
\begin{teile}
\steil Eine Wandergruppe geht mit 4 km/h. Sie ist um 9:50 Uhr an einer Hütte, macht dort 30 Minuten Pause und geht dann 3\,000 m bis zum See. Wann kommt sie am See an? \hfill (P10 2014 OS) \\ um \leerfeld Uhr
\steil Ein Radfahrer fährt mit 18 km/h. Er startet um 14:40 Uhr, macht unterwegs 20 Minuten Pause und fährt insgesamt 9\,000 m. Wann kommt er an? \hfill (P10 2014 OS) \\ um \leerfeld Uhr
\steil Die Erde und der Mars sind an einem Tag ca. $2{,}25 \cdot 10^{8}$ km voneinander entfernt. Ein Funksignal legt ca. 300\,000 km in einer Sekunde zurück. Wie viele Sekunden braucht ein Signal von der Erde zum Mars? Wie viele Minuten sind das? \hfill (P10 2015 OS) \\ \leerfeld[s] ; \leerfeld[min]
\steil Das Licht der Sonne legt bis zur Erde ca. $1{,}5 \cdot 10^{8}$ km zurück. Licht schafft ca. 300\,000 km in einer Sekunde. Wie viele Sekunden braucht das Licht? Wie viele Minuten sind das? Runde auf eine Stelle nach dem Komma. \hfill (P10 2015 OS) \\ \leerfeld[s] ; \leerfeld[min]
\teil Eine Seilbahn fährt eine 630 m lange Strecke mit durchschnittlich 3,5 m/s. Wie viele Minuten dauert die Fahrt? \hfill (P10 2024 OS) \\ \leerfeld[min]
\teil Ein Förderband transportiert Koffer 420 m weit, im Schnitt mit 1,4 m/s. Wie viele Minuten ist ein Koffer unterwegs? \hfill (P10 2024 OS) \\ \leerfeld[min]
\end{teile}
\end{aufgabe}
